\section{Reconstruction and the map to a complex}\label{sec:reconstruction}

All distances in this section are measured in the original metric $g$;
write $d_0$ for that distance.  Put $D=WR_0$, as in
\cref{rounds:wide}.  The centers $x_i$ form a maximal $D$-separated set,
and the color pairs are
\[
 A_i=\overline B_0(x_i,D),\qquad
 B_i=M\setminus B_0(x_i,2D).
\]
The $A_i$ cover $M$.  Every family of balls with these centers and a
fixed finite radius is locally finite: a compact set has a compact
neighborhood of that radius by completeness, and a compact metric
space contains only finitely many $D$-separated points.  The family is
countable, since $M$ is second countable.  If some $B_i$ is empty,
$\diam_{g}M\le 4D$ and a constant map proves the theorem.  We therefore
assume that all $B_i$ are nonempty.

\subsection{A limit on the original manifold}

The analytic packets have infinite product ends.  Their finite
reconstructions in \cref{packet:reconstruction} are manifolds
homeomorphic to $M$; only finite intervals of the product ends are used
in these reconstructions.  We now choose those homeomorphisms
coherently.  This is the point at which local finiteness of the color
supports is essential.

\begin{proposition}[Reconstructed partitions]\label{topology:partitions}
The compatible cuts of \cref{countable:construction} give a continuous
proper map $q:M\to M$, and for every color $i$ closed sets $P_i,N_i,W_i$
such that
\begin{align}
 &d_0(q(x),x)<1,\label{topology:displacement}\\
 &P_i\cup N_i=M,\qquad P_i\cap N_i=W_i,
   \qquad W_i\subset q^{-1}\bigl(B_0(x_i,2D)\bigr),
   \label{topology:sides}\\
 &q^{-1}(A_i)\subset P_i\setminus W_i,
   \qquad q^{-1}(B_i)\subset N_i\setminus W_i.
   \label{topology:forced}
\end{align}
The wall family $(W_i)$ is locally finite, and no $n-1$ distinct walls
have a common point.
\end{proposition}

\begin{proof}
Choose positive numbers $\eps_0,\eps_1,\ldots$ with
$\sum_{r\ge0}\eps_r<1$.  Here $\eps_0$ is assigned to the single radial
stage; the other numbers are assigned to successive color stages.
In this proof only, index the post-radial reconstruction by $r=0$.
For a finite reconstruction write
\[
 q_r:X_r\longrightarrow M,\qquad
 H_r:X_r\longrightarrow M,\qquad
 Q_r=q_r\circ H_r^{-1}.
\]
The map $q_r$ is the proper collapse of
\cref{packet:reconstruction}, whereas $H_r$ is a homeomorphism to a
fixed copy of the original manifold.  We maintain
\begin{equation}\label{topology:finite-error}
 d_0(Q_r(x),x)\le E_r,
 \qquad E_r\le\sum_{j=0}^{r}\eps_j<1.
\end{equation}
Walls and their two closed sides will also be viewed on this fixed copy
using $H_r$.

For the radial stage, the cut has a locally finite family of compact
components.  Choose pairwise disjoint, locally finite collars of these
components on which the collapse varies by less than $\eps_0$ in each
normal segment.  Compactness of each component and continuity of the
collapse permit the collars to be this small, independently from one
component to another.  Compress the finite inserted cylinders into
these collars.  Since their supports are disjoint, the displacement
bound is $\eps_0$, rather than a sum over components.  This constructs
$H_0$ and proves \eqref{topology:finite-error} initially.  The finite
length of an inserted cylinder may be arbitrarily large; its length
does not impose a lower bound on the width of the collar into which it
is compressed.

We describe a subsequent color stage $i$.  First raise any old finite
horizons needed for the new compatible defining functions.  Denote the
resulting reconstruction by $\widehat X_r$.  The horizon
reparameterization
\[
 \Phi_r:\widehat X_r\longrightarrow X_r
\]
preserves collapse and every old wall and signed side.  Indeed, on an
old product end the collapse is independent of the normal coordinate,
and an increasing change of that coordinate preserves its sign and its
zero section.  Compatible products of these changes agree at corners.
Thus we may set
\begin{equation}\label{topology:horizon}
 \widehat q_r=q_r\circ \Phi_r,\qquad
 \widehat H_r=H_r\circ \Phi_r,
 \qquad
 \widehat q_r\circ\widehat H_r^{-1}=Q_r.
\end{equation}
In particular, increasing a horizon contributes no displacement on the
fixed copy.

The new wall lies over $B_0(x_i,2D)$.  Choose its collar strip $S_i$ in
$\widehat X_r$ so that
\begin{equation}\label{topology:strip}
 S_i\subset\widehat q_r^{-1}\bigl(B_0(x_i,2D+1)\bigr),
\end{equation}
and so that collapse oscillation along every segment used by the
compression is less than the error assigned to this stage.  These
choices are possible because the relevant wall components in a finite
reconstruction are compact and locally finite.  More specifically,
their union lies in the inverse image of the compact ball
$\overline B_0(x_i,2D)$ under the proper finite collapse; hence only
finitely many components are involved in this color stage.

The strips can be chosen to preserve every old stratum, including the
radial strata.  To check this assertion, work near a point of a new wall
on an old stratum, and denote the new defining function by $u$.
The horizons were chosen beyond the product thresholds, so the terminal
collars glued at each old hit wall lie in the product range.  In the
resulting product chart, $u$ is pulled back from the corresponding face
and is independent of the old signed normal coordinates.  Its restriction
to that face is regular at zero.  In the corresponding product chart choose a vector field
$V$ in the face directions with $V(u)=1$, and with zero component in
every old normal direction.  The chart may avoid the old walls not
passing through the point.  A partition of unity combines these fields:
the linear identity $V(u)=1$ and tangency to every old wall both survive.
Its short flow supplies a collar preserving the old walls and their
sides.  On a fresh cylinder, each old stratum is the product of its
intersection with the new wall and the cylinder interval.  The
compression uses the same increasing normal reparameterization on
all these products.  It consequently preserves every old signed side
as well as all their intersections.

Let
\[
 c_r:X_{r+1}\longrightarrow\widehat X_r
 \quad\hbox{and}\quad
 h_r:X_{r+1}\longrightarrow\widehat X_r
\]
be, respectively, the new cylinder collapse and the homeomorphism that
compresses those cylinders into the chosen strips.  They agree off
the inverse image of $S_i$. Figure~\ref{fig:topology-compression} shows
the two maps schematically in a normal coordinate: $c_r$ collapses the
inserted interval, while $h_r$ compresses it into the chosen old collar.
Define
\begin{equation}\label{topology:recursion}
 q_{r+1}=\widehat q_r\circ c_r,
 \qquad H_{r+1}=\widehat H_r\circ h_r,
 \qquad
 Q_{r+1}=\widehat q_r\circ c_r\circ h_r^{-1}
                         \circ\widehat H_r^{-1}.
\end{equation}
The oscillation choice gives
$d_0(Q_{r+1}(x),Q_r(x))<\eps_{r+1}$ wherever they differ, proving
\eqref{topology:finite-error} at the next stage.  Their difference is
supported in $\widehat H_r(S_i)$.  If
$x=\widehat H_r(z)$ belongs to this set, then
\eqref{topology:horizon}, \eqref{topology:strip}, and
\eqref{topology:finite-error} imply
\begin{equation}\label{topology:fixed-support}
 d_0(x,x_i)
 \le d_0(x,Q_r(x))+d_0(\widehat q_r(z),x_i)
 <1+(2D+1)=2D+2.
\end{equation}
This support radius is independent of the stage.

\begin{figure}[!b]
\centering
\begin{tikzpicture}[x=1.25cm,y=1.05cm,>=Latex,font=\small]
 \fill[blue!6] (-3,-1) rectangle (3,1);
 \fill[red!8] (-1,-2.25) rectangle (1,2.35);
 \draw[->] (-3.2,0)--(3.5,0) node[right] {$s$};
 \draw[->] (0,-2.3)--(0,2.65) node[above] {$t$};
 \draw[densely dashed,gray] (-3,1)--(3,1);
 \draw[densely dashed,gray] (-3,-1)--(3,-1);
 \draw[very thick,purple!75!black]
  (-3,-2)--(-1,0)--(1,0)--(3,2);
 \draw[thick,orange!85!black]
  (-3,-2)--(-2,-1)--(-1,-.35)--(0,0)--(1,.35)--(2,1)--(3,2);
 \draw (-1,.07)--(-1,-.07) node[below=4pt,fill=white,inner sep=1pt] {$-L$};
 \draw (1,.07)--(1,-.07) node[below=4pt] {$L$};
 \node[left,fill=white,inner sep=1pt] at (-.08,1) {$\delta$};
 \node[left,fill=white,inner sep=1pt] at (-.08,-1) {$-\delta$};
 \node[align=left,anchor=west,text=purple!75!black] at (3.55,1.25)
  {collapse $c_r$};
 \node[align=left,anchor=west,text=orange!85!black] at (3.55,.55)
  {compression $h_r$};
 \node[align=left,anchor=west] at (3.55,-.55)
  {$c_r=h_r$ outside the strip};
 \node[align=left,anchor=west] at (3.55,-1.35)
  {old strata follow\\the same reparameterization};
 \node[below,align=center] at (0,-2.4)
  {finite inserted cylinder: $-L\le s\le L$};
\end{tikzpicture}
\caption{The two maps used in a finite reconstruction, shown schematically
in local normal coordinates: $s$ on the reconstructed domain $X_{r+1}$
and $t$ on the old target $\widehat X_r$. The target strip $|t|<\delta$
represents the chosen collar $S_i$.
The collapse is constant on the inserted cylinder;
the homeomorphism squeezes it into an arbitrarily narrow old collar.
Where the maps differ, both images lie in the target strip.  Infinite analytic ends
are truncated before this topological operation.}
\label{fig:topology-compression}
\label{packet:collar-figure}
\end{figure}

The balls in \eqref{topology:fixed-support} form a locally finite
family.  Consequently, on every relatively compact neighborhood in
the fixed copy, the maps $Q_r$ are exactly unchanged after finitely
many stages.  Their locally stationary limit $q$ is continuous and
satisfies \eqref{topology:displacement}.  It is proper: the inverse
image of a compact set is a closed subset of its compact closed
$1$-neighborhood in $(M,d_0)$.

At each finite stage the signed sides are the closed reconstructed
sides, with their intersection equal to the central wall.  The
stratum-preserving identifications just constructed preserve every
previously installed side and wall exactly.  Newly installed color
walls are confined to the fixed balls in
\eqref{topology:fixed-support}.  Thus all the partition data are locally
eventually unchanged.  Their final closed sets satisfy
\eqref{topology:sides} and \eqref{topology:forced}, since these are the
finite-stage forced-sign properties of
\cref{countable:construction}, and both collapse and statuses are
stationary near each point.  The wall family is locally finite by the
same support bound.  Finally, a point on $n-1$ distinct final color
walls would, on a neighborhood where these data have stabilized, lie
on the same $n-1$ finite-stage walls.  This contradicts
\cref{countable:construction}.
\end{proof}

\subsection{Thickening the walls without increasing their multiplicity}

Put $k=n-2$.  We need open neighborhoods of the whole walls, rather
than independently chosen neighborhoods of their individual compact
pieces.

\begin{lemma}[Simultaneous wall neighborhoods]\label{topology:thicken}
There are locally finite open sets $T_i\supset W_i$ such that
\begin{equation}\label{topology:thickening}
 T_i\subset q^{-1}\bigl(B_0(x_i,2D)\bigr)
                  \setminus q^{-1}(A_i\cup B_i),
 \qquad
 \#\{i:x\in T_i\}\le k\quad(x\in M).
\end{equation}
\end{lemma}

\begin{proof}
Set $O_i=q^{-1}(B_0(x_i,2D))$ and
$F_i=q^{-1}(A_i\cup B_i)$.  By
\cref{topology:partitions}, $W_i$ is disjoint from the closed set $F_i$
and lies in the open set $O_i$.  At a point $x$, local finiteness and
closedness of the walls give an open neighborhood meeting only walls
that contain $x$.  Shrink it to a neighborhood $U_x$ contained in
$O_i\setminus F_i$ for every such wall.  Therefore $U_x$ meets at most
$k$ walls, and
\begin{equation}\label{topology:cover-property}
 U_x\cap W_i\ne\varnothing\quad\Longrightarrow\quad
 U_x\subset O_i\setminus F_i.
\end{equation}

Choose a locally finite open star refinement $\mathcal V$ of this
cover: the union of the members of $\mathcal V$ meeting any fixed
member is contained in some $U_x$.  For completeness, such a refinement
can be obtained directly with the metric $d_0$.  The function
\[
 \rho(x)=\sup_U\min\{1,d_0(x,M\setminus U)\},
 \qquad U\in\{U_x:x\in M\},
\]
is positive and $1$-Lipschitz, with distance to the empty set truncated
to $1$.  The balls $V_x=B_0(x,\rho(x)/8)$ star-refine the original
cover.  Indeed, if $V_y$ meets $V_x$, the Lipschitz inequality gives
$d_0(x,y)<2\rho(x)/7$, and hence
$V_y\subset B_0(x,25\rho(x)/56)$.  The latter ball is contained in
$B_0(x,\rho(x)/2)$, which lies in a member of the original cover.
Fix $o\in M$.  Obtain a locally finite open refinement of $(V_x)$ by covering each
compact annulus
$\overline B_0(o,m)\setminus B_0(o,m-1)$ with finitely many open subsets
of members of $(V_x)$ lying in
$\{m-2<d_0(o,\cdot)<m+1\}$, for $m\ge1$.
These sets cover $M$ and are locally finite.  A refinement of a star
refinement is still a star refinement of the original cover, which
proves the assertion about $\mathcal V$.

Define
\[
 T_i=\bigcup\{V\in\mathcal V:V\cap W_i\ne\varnothing\}.
\]
It contains $W_i$, and \eqref{topology:cover-property} implies the
containment in \eqref{topology:thickening}.  If $z$ belongs to several
$T_i$, choose $V_0\in\mathcal V$ containing $z$.  For each such $i$
there is a $V_i\in\mathcal V$ containing $z$ and meeting $W_i$.
All the $V_i$ meet $V_0$, so their union lies in a single $U_x$.
That $U_x$ meets each of these walls, and therefore there are at most
$k$ of them.  Finally, $(T_i)$ is locally finite because $T_i\subset O_i$
and the family $(O_i)$ is locally finite.
\end{proof}

\subsection{Coordinates, triangulation, and entire fibers}

For each $i$ the closed sets
\[
 C_i^+=P_i\setminus T_i,\qquad C_i^-=N_i\setminus T_i
\]
are disjoint and cover $M\setminus T_i$.  For a closed set $C$ write
$\delta_C(x)=\min\{1,d_0(x,C)\}$, with $\delta_\varnothing=1$, and put
\begin{equation}\label{topology:coordinates}
 t_i(x)=\frac{\delta_{C_i^-}(x)}
                    {\delta_{C_i^-}(x)+\delta_{C_i^+}(x)}.
\end{equation}
The denominator is positive because the two closed sets are disjoint.
These continuous coordinates have the following properties:
\begin{enumerate}[label=(\roman*)]
 \item $t_i=1$ on $q^{-1}(A_i)$ and $t_i=0$ on $q^{-1}(B_i)$;
 \item $0<t_i(x)<1$ implies $x\in T_i$;
 \item $\supp t_i\subset q^{-1}(\overline B_0(x_i,2D))$,
       so the coordinate supports are locally finite;
 \item at every point at least one coordinate equals $1$, and at most
       $k$ coordinates lie strictly between $0$ and $1$.
\end{enumerate}
Here the last assertion uses the inner-ball cover and
\cref{topology:thicken}.

Consider the cubical set of finitely supported coordinate vectors
described by these properties.  Its cubes are
\[
 Q(S,A)=\{t:t_i=1\ (i\in S),\ 0\le t_i\le1\ (i\in A),
                    \ t_i=0\ (i\notin S\cup A)\},
\]
where $S,A$ are disjoint finite sets of labels, $S\ne\varnothing$, and
$|A|\le k$.  Triangulate each cube by the staircase triangulation: for
every ordering $a_1,\ldots,a_j$ of $A$, take the simplex with vertices
\begin{equation}\label{topology:staircase}
 \mathbf1_S,\quad \mathbf1_{S\cup\{a_1\}},\quad\ldots,\quad
 \mathbf1_{S\cup A}.
\end{equation}
The triangulations agree on faces: setting a free coordinate to zero
removes it, while setting it to one moves it to $S$.  They therefore
define a simplicial complex $K$ of dimension at most $k$.

We use the usual polyhedral topology on $|K|$.  To verify continuity
without placing a product topology on the infinite coordinate space,
fix a point of $M$ and a neighborhood meeting only finitely many
coordinate supports.  On that neighborhood all other coordinates
are identically zero.  The vector $(t_i)$ thus takes values in the
finite cubical complex using this finite set of labels, and is a
continuous map to that complex.  Its staircase triangulation is a
finite geometric simplicial complex, whose inclusion in $|K|$ is
continuous.  The coordinate map consequently defines a continuous map
\begin{equation}\label{topology:final-map}
 F:M\longrightarrow |K|,\qquad \dim K\le n-2.
\end{equation}
The triangulation retains the entire coordinate vector: the affine
map from each simplex in \eqref{topology:staircase} to its coordinate
simplex agrees on faces and is injective on their union.  One can see
injectivity by restricting any two finite supports to their finite
union; there it is the usual triangulation of a cubical complex.
Thus $F(x)=F(y)$ implies $t_i(x)=t_i(y)$ for every $i$.

Choose, for any two points of the same fiber, a coordinate that equals
$1$ at the first point.  It also equals $1$ at the second.  By (i), both
collapse images lie in the same open outer ball, and hence
\[
 d_0(x,y)
 \le d_0(x,q(x))+d_0(q(x),x_i)
       +d_0(x_i,q(y))+d_0(q(y),y)
 <4D+2.
\]
This applies to every pair of points in the fiber, regardless of its
number of components.  Consequently
\begin{equation}\label{topology:fiber-bound}
 \sup_{z\in|K|}\diam_{g}F^{-1}(z)\le4D+2=4WR_0+2.
\end{equation}
Together with \cref{rounds:wide,countable:construction}, this proves
\cref{thm:main}.
